% GENERATED FILE. Edit canonical lessons, then run npm run generate:books.
% canonical-source-hash: fnv1a64:55fbb93d4073a112
% canonical-lessons: IT-C533-apparecchiare, IT-C533-sparecchiare, IT-C533-tovaglia, IT-C533-tovagliolo, IT-C533-vassoio

\chapter{A tavola}
\label{ch:it-a-tavola}
\hlchaptermodality{533}

\begin{chapteropening}
\textbf{By the end of this chapter:} \emph{I can talk about laying and clearing a table for a meal.}

Complete the last lesson of chapter 533: I can talk about laying and clearing a table for a meal.
\end{chapteropening}

\section[apparecchiare]{apparecchiare — to lay (a table)}

\label{lesson:IT-C533-apparecchiare}



\begin{quote}
Before the new one: say the Italian for detergent, then the Italian for vacuum cleaner.
\end{quote}

\subsection*{You'll want to know: apparecchiare}
\textbf{apparecchiare} — \textquotedblleft{}to lay (a table)\textquotedblright{}.

\begin{quote}
Apparecchio la tavola per sei persone.
\end{quote}

\begin{quote}
\emph{I lay the dinner table for six people.}
\end{quote}

\emph{Apparecchiare la tavola}.

\begin{cousinweb}
\emph{apparecchiare} is usually traced to Vulgar Latin \emph{appariculāre} \textquotedblleft{}to prepare\textquotedblright{}, from Latin \emph{apparāre}, built on \emph{parāre} \textquotedblleft{}to make ready\textquotedblright{} (Lewis \& Short).
\end{cousinweb}

\begin{grammarlens}[title={What you've built}]
One more word for telling, explaining and describing at B1.
\end{grammarlens}

\subsection*{Your turn}
\begin{itemize}
\raggedright
  \item \emph{Say it:} \emph{apparecchiare}
  \item \emph{Say it:} \emph{apparecchiare}, once more
  \item \emph{Say it:} a sentence of your own with \emph{apparecchiare}
  \item \emph{Recall:} say the Italian for detergent, then the Italian for vacuum cleaner, then say \emph{apparecchiare} again
\end{itemize}

\begin{tcolorbox}[breakable,colback=teal!4,colframe=teal!35!black,title={Before you move on}]
What does apparecchiare mean? (\textquotedblleft{}To lay (a table)\textquotedblright{}.) Say it once more.
\end{tcolorbox}
\section[sparecchiare]{sparecchiare — to clear (a table)}

\label{lesson:IT-C533-sparecchiare}



\begin{quote}
Before the new one: say the Italian for vacuum cleaner, then the Italian for to lay.
\end{quote}

\subsection*{You'll want to know: sparecchiare}
\textbf{sparecchiare} — \textquotedblleft{}to clear (a table)\textquotedblright{}.

\begin{quote}
Chi sparecchia stasera?
\end{quote}

\begin{quote}
\emph{Who is clearing the dinner table tonight?}
\end{quote}

The opposite of \emph{apparecchiare}.

\begin{cousinweb}
\emph{sparecchiare} is built from \emph{s-} \textquotedblleft{}un-\textquotedblright{} and \emph{apparecchiare} \textquotedblleft{}to lay a table\textquotedblright{}, which goes back to Latin \emph{apparāre} \textquotedblleft{}to prepare\textquotedblright{} (Lewis \& Short).
\end{cousinweb}

\begin{grammarlens}[title={What you've built}]
One more word for telling, explaining and describing at B1.
\end{grammarlens}

\subsection*{Your turn}
\begin{itemize}
\raggedright
  \item \emph{Say it:} \emph{sparecchiare}
  \item \emph{Say it:} \emph{sparecchiare}, once more
  \item \emph{Say it:} a sentence of your own with \emph{sparecchiare}
  \item \emph{Recall:} say the Italian for vacuum cleaner, then the Italian for to lay, then say \emph{sparecchiare} again
\end{itemize}

\begin{tcolorbox}[breakable,colback=teal!4,colframe=teal!35!black,title={Before you move on}]
What does sparecchiare mean? (\textquotedblleft{}To clear (a table)\textquotedblright{}.) Say it once more.
\end{tcolorbox}
\section[la tovaglia]{la tovaglia — tablecloth}

\label{lesson:IT-C533-tovaglia}



\begin{quote}
Before the new one: say the Italian for to lay, then the Italian for to clear.
\end{quote}

\subsection*{You'll want to know: la tovaglia}
\textbf{la tovaglia} — \textquotedblleft{}tablecloth\textquotedblright{}.

\begin{quote}
Per la festa mettiamo la tovaglia bianca.
\end{quote}

\begin{quote}
\emph{For the party we put on the white tablecloth.}
\end{quote}

\emph{Una tovaglia a quadri} — a checked tablecloth.

\begin{cousinweb}
\emph{tovaglia} comes from a Germanic word, Frankish \emph{thwahillja} \textquotedblleft{}a cloth\textquotedblright{} (EtymDB, under English \emph{towel}), the same source as English \emph{towel}.
\end{cousinweb}

\begin{grammarlens}[title={What you've built}]
One more word for telling, explaining and describing at B1.
\end{grammarlens}

\subsection*{Your turn}
\begin{itemize}
\raggedright
  \item \emph{Say it:} \emph{la tovaglia}
  \item \emph{Say it:} \emph{la tovaglia}, once more
  \item \emph{Say it:} a sentence of your own with \emph{la tovaglia}
  \item \emph{Recall:} say the Italian for to lay, then the Italian for to clear, then say \emph{la tovaglia} again
\end{itemize}

\begin{tcolorbox}[breakable,colback=teal!4,colframe=teal!35!black,title={Before you move on}]
What does la tovaglia mean? (\textquotedblleft{}Tablecloth\textquotedblright{}.) Say it once more.
\end{tcolorbox}
\section[il tovagliolo]{il tovagliolo — napkin}

\label{lesson:IT-C533-tovagliolo}



\begin{quote}
Before the new one: say the Italian for to clear, then the Italian for tablecloth.
\end{quote}

\subsection*{You'll want to know: il tovagliolo}
\textbf{il tovagliolo} — \textquotedblleft{}napkin\textquotedblright{}.

\begin{quote}
Mi passi un tovagliolo, per favore?
\end{quote}

\begin{quote}
\emph{Can you pass me a napkin, please?}
\end{quote}

Built on \emph{tovaglia}.

\begin{cousinweb}
\emph{tovagliolo} is a little form of \emph{tovaglia} \textquotedblleft{}tablecloth\textquotedblright{}, which comes from Frankish \emph{thwahillja} \textquotedblleft{}a cloth\textquotedblright{} (EtymDB).
\end{cousinweb}

\begin{grammarlens}[title={What you've built}]
One more word for telling, explaining and describing at B1.
\end{grammarlens}

\subsection*{Your turn}
\begin{itemize}
\raggedright
  \item \emph{Say it:} \emph{il tovagliolo}
  \item \emph{Say it:} \emph{il tovagliolo}, once more
  \item \emph{Say it:} a sentence of your own with \emph{il tovagliolo}
  \item \emph{Recall:} say the Italian for to clear, then the Italian for tablecloth, then say \emph{il tovagliolo} again
\end{itemize}

\begin{tcolorbox}[breakable,colback=teal!4,colframe=teal!35!black,title={Before you move on}]
What does il tovagliolo mean? (\textquotedblleft{}Napkin\textquotedblright{}.) Say it once more.
\end{tcolorbox}
\section[il vassoio]{il vassoio — tray}

\label{lesson:IT-C533-vassoio}



\begin{quote}
Before the new one: say the Italian for tablecloth, then the Italian for napkin.
\end{quote}

\subsection*{You'll want to know: il vassoio}
\textbf{il vassoio} — \textquotedblleft{}tray\textquotedblright{}.

\begin{quote}
Porto il caffè sul vassoio.
\end{quote}

\begin{quote}
\emph{I bring the coffee on the tray.}
\end{quote}

\emph{Un vassoio di pasticcini} — a tray of pastries.

\begin{cousinweb}
\emph{vassoio} is usually traced to Vulgar Latin \emph{missōrium} \textquotedblleft{}a serving dish\textquotedblright{} (EtymDB), from \emph{mittere} \textquotedblleft{}to send\textquotedblright{}; the details are not settled.
\end{cousinweb}

\begin{grammarlens}[title={What you've built}]
One more word for telling, explaining and describing at B1.
\end{grammarlens}

\subsection*{Your turn}
\begin{itemize}
\raggedright
  \item \emph{Say it:} \emph{il vassoio}
  \item \emph{Say it:} \emph{il vassoio}, once more
  \item \emph{Say it:} a sentence of your own with \emph{il vassoio}
  \item \emph{Recall:} say the Italian for tablecloth, then the Italian for napkin, then say \emph{il vassoio} again
\end{itemize}

\begin{tcolorbox}[breakable,colback=teal!4,colframe=teal!35!black,title={Before you move on}]
What does il vassoio mean? (\textquotedblleft{}Tray\textquotedblright{}.) Say it once more.
\end{tcolorbox}
